% ==============================================================================
% =================================== Aufgabe 4 =================================
% ==============================================================================

\chapter{Aufgabe}
\label{chap:4 Aufgabe}
\thispagestyle{fancy}
\section{Einleitungskapitel} 
\begin{enumerate}[label=\thechapter.\arabic*), leftmargin=*]
    \item \textbf{Einleitung}
    \begin{enumerate}[label=\thechapter.\arabic{enumi}.\arabic*), leftmargin=*]
        \item \textbf{Relevanz des Themas und Motivation} \par In einer globalisierten Arbeitswelt, die durch zunehmende Volatilität, Unsicherheit, Komplexität und Mehrdeutigkeit (\textit{\acs{VUCA}}) geprägt ist, stoßen traditionelle, starre Linienorganisationen vermehrt an ihre Grenzen. Um schnell und flexibel auf Marktanforderungen reagieren zu können, vollziehen insbesondere Unternehmen der \acs{IT}-Branche eine „\textbf{\textit{agile Transformation}}“. Während das klassische Projektmanagement auf detaillierter Vorabplanung und zentraler Steuerung basiert, setzt das agile Paradigma auf Flexibilität und iterative Prozesse. \par Die vorliegende Arbeit entspringt der Motivation, den menschlichen Faktor in diesem Transformationsprozess zu untersuchen. Studien belegen, dass der Erfolg eines Projekts maßgeblich von den sogenannten „\textbf{\textit{Soft Facts}}“ – wie Führung, Kommunikation und Motivation – abhängt, während strukturelle „\textbf{\textit{Hard Facts}}“ lediglich unterstützend wirken. Ein zentraler Aspekt ist dabei der Rollenwechsel der Führungskraft: Weg vom klassischen Projektleiter, der als „\textbf{\textit{Unternehmer des Projekts}}“ agiert, hin zum Scrum Master, der als „\textbf{\textit{Servant Leader}}“ Hindernisse beseitigt und die Selbstorganisation des Teams coacht. Da Softwareprojekte nicht nur technische, sondern primär verhaltensbedingte Herausforderungen darstellen, ist es von hoher wissenschaftlicher und praktischer Relevanz, zu klären, wie sich dieser Rollenwandel auf die intrinsische Motivation der Teams auswirkt.
        \item \textbf{Zielsetzung der Arbeit und Arbeitshypothese} \par Das primäre Ziel dieser Arbeit ist es, die Veränderung der Motivationsstrukturen innerhalb von Softwareentwicklungsteams zu analysieren, wenn die hierarchische Kontrolle durch agile Selbststeuerung ersetzt wird. Dabei wird untersucht, ob die geänderte Führungsdynamik des Scrum Masters ein Umfeld schafft, das die Eigenverantwortung und damit die Arbeitszufriedenheit nachhaltig steigert. \par Als forschungsleitende \textbf{Arbeitshypothese} wird aufgestellt:
            \cleardoublepage
            \begin{itemize}
                \item \textbf{Hypothese:} Der Übergang von der zentralen Steuerung durch einen klassischen Projektleiter hin zur unterstützenden Moderation durch einen Scrum Master korreliert positiv mit der intrinsischen Motivation des Projektteams.
            \end{itemize}
            Das Ziel der Arbeit ist erreicht, wenn der Zusammenhang zwischen diesen beiden Variablen (\textit{Rollenmodell als Variable A und Motivation als Variable B}) eindeutig identifiziert und bewertet werden konnte.
        \item \textbf{Wissenschaftliche Vorgehensweise} \par Um eine robuste und valide Beantwortung der Forschungsfrage zu gewährleisten, wird ein methodischer Triangulationsansatz gewählt. Das Vorgehen gliedert sich in drei Säulen:
            \begin{enumerate}[label=\thechapter.\arabic{enumi}.\arabic{enumii}.\arabic*), leftmargin=*]
                \item \textbf{Systematisches Review:} Zunächst erfolgt eine systematische Auswertung aktueller Fachliteratur und empirischer Studien der letzten sechs Jahre, um den aktuellen Forschungsstand zu Rollenmodellen und Team-Motivation zu aggregieren.
                \item \textbf{Quantitative Befragung:} Mithilfe eines standardisierten Online-Fragebogens werden Daten von Softwareentwicklern erhoben. Die statistische Analyse fokussiert sich auf die Berechnung von Korrelationskoeffizienten, um den Zusammenhang zwischen Führungsstil und Motivation quantitativ zu belegen.
                \item \textbf{Qualitative Experteninterviews:} Ergänzend werden leitfadengestützte Interviews mit erfahrenen Scrum Mastern geführt. Die Auswertung erfolgt mittels der qualitativen Inhaltsanalyse nach Mayring, um tiefergehende Einblicke in individuelle Motivationsfaktoren zu gewinnen.
            \end{enumerate}
        \item \textbf{Aufbau der Arbeit} \par Die Arbeit ist in sechs Kapitel gegliedert, die einer logischen wissenschaftlichen Struktur folgen:
            \begin{itemize}
                \item \textbf{Kapitel 2 (\textit{Theoretische Fundierung}):} Hier werden die konzeptionellen Grundlagen des Projektmanagements erarbeitet. Es erfolgt eine detaillierte Gegenüberstellung der Rollenprofile (\textit{Projektleiter \acs{vs.} Scrum Master}) sowie eine Einordnung relevanter Motivationstheorien.
                \item \textbf{Kapitel 3 (\textit{Methodik}):} In diesem Kapitel wird das Forschungsdesign umfassend begründet. Die Wahl der Erhebungsinstrumente (\textit{Fragebogen und Interviewleitfaden}) sowie der Stichprobenaufbau werden dargelegt.
                \item \textbf{Kapitel 4 (\textit{Empirische Untersuchung und Ergebnisanalyse}):} Hier werden die im Feld erhobenen Daten präsentiert. Die statistischen Korrelationen und die zentralen Aussagen aus den Interviews werden sachlich dargestellt.
                \item \textbf{Kapitel 5 (\textit{Diskussion und Handlungsempfehlungen}):} Die Ergebnisse werden vor dem Hintergrund der Hypothese interpretiert. Es werden Limitationen der Studie aufgezeigt und konkrete Empfehlungen für das \acs{IT}-Management abgeleitet.
                \item \textbf{Kapitel 6 (\textit{Schlussfolgerung}):} Die Arbeit schließt mit einer Zusammenfassung der Kernerkenntnisse und einem Ausblick auf zukünftige Forschungsbedarfe im Bereich der agilen Führung.
            \end{itemize}
    \end{enumerate}
\end{enumerate}

% ==============================================================================
% =================================== Aufgabe 4 =================================
% ==============================================================================